



% \subsection{Memory Benchmarks}
% \label{subsec:bench_memory}




% Evaluating memory in LLM-based agents has historically depended on a fragmented ecosystem of datasets, each narrowly targeting a specific ability such as long-context retrieval~\cite{hsieh2024ruler, wu2024longmemeval, zhang2024bench}, dialogue and speaker-state tracking~\cite{li2018towards, liu2021benchmarking}, preference or recommendation learning~\cite{casanueva2020efficient, larson2019evaluation, he2023large}, or counterfactual knowledge editing and revision~\cite{zhong2023mquake}. While individually valuable, these datasets lack a unified framework for understanding how agents retain, integrate, synthesize, or selectively forget information across multi-turn interactions. To address this gap, we adopt \emph{MemoryAgentBench}~\cite{hu2025memoryagentbench}, a consolidated benchmark constructed to measure four foundational memory competencies in LLM-driven agents: \textit{Accurate Retrieval (AR)}, \textit{Test-Time Learning (TTL)}, \textit{Long-Range Understanding (LRU)}, and \textit{Selective Forgetting (SF)}, thereby providing a principled and comprehensive diagnostic lens on memory-centric behavior. These competencies correspond to essential memory operations. AR probes precise recall and multi-hop retrieval, TTL captures incremental concept and preference acquisition, LRU stresses global abstraction over long-form content, and SF targets controlled revision of conflicting or outdated information.


% MemoryAgentBench~\footnote{\url{https://huggingface.co/datasets/ai-hyz/MemoryAgentBench}} organizes its benchmark suite into four splits, each aligned with one competency, and each containing multiple subtasks adapted from diverse domains. The \textit{AR} split combines four datasets: \textit{SH-QA}, \textit{MH-QA}, \textit{LME(S*)}, and \textit{EventQA}, which evaluate single-hop factual recall, multi-hop reasoning, long-dialogue memory tracking, and temporal event understanding. SH-QA and MH-QA are adapted from long-context QA resources such as RULER~\cite{hsieh2024ruler} and LongMemEval~\cite{wu2024longmemeval}, while LME(S*) derives from dialogue-based long-memory evaluations, and EventQA extends temporal reasoning benchmarks from $\infty$-Bench~\cite{zhang2024bench}. The \textit{TTL} split evaluates rapid in-session learning through two subtasks: \textit{MCC}, adapted from intent-classification datasets~\cite{casanueva2020efficient, larson2019evaluation}, and \textit{Recom.}, derived from recommendation and preference-learning corpora~\cite{li2002learning, liu2021benchmarking, li2018towards, he2023large}. These tasks assess whether an agent can incorporate newly introduced concepts, categories, or user preferences without requiring retraining. The \textit{LRU} split comprises \textit{Summ.} and \textit{DetQA}, which respectively test long-range abstraction and cross-document global reasoning. Summ. originates from the long-context summarization tasks in $\infty$-Bench~\cite{zhang2024bench}, whereas DetQA~\cite{xu2024detectiveqa} introduces narrative-spanning inference in detective-style scenarios where relevant evidence is dispersed across distant portions of the text. Finally, the \textit{SF} split introduces two fact-consolidation subtasks: \textit{FC-SH} and \textit{FC-MH}, constructed from counterfactual edits in MQUAKE~\cite{zhong2023mquake}. These tasks probe whether an agent can correctly incorporate updated information while overriding earlier contradictions, capturing a controlled form of selective forgetting necessary for maintaining coherent long-term memory representations.










% \input{figures/MemoryInterface}
















% Beyond consolidating datasets, MemoryAgentBench prescribes a uniform agent interface that enables consistent comparison across heterogeneous frameworks; however, realizing such comparisons in practice requires substantial engineering effort. In this work, we explicitly implement the MemoryAgentBench interface across all evaluated agent frameworks, standardizing memory initialization, incremental ingestion, and query execution so that observed performance differences reflect architectural and memory design choices rather than inconsistencies in experimental setup. As shown in Listing~\ref{lst:memoryagent-interface}, every agent implements three core operations: \texttt{reset}, which initializes or clears all memory stores; \texttt{ingest}, which processes incrementally provided context segments; and \texttt{query}, which answers memory-intensive questions using the accumulated history. Each benchmark session provides a context document and a set of questions, and we follow the official specification by dividing each context into fixed-size chunks of up to 4,096 tokens with a 200-token overlap and delivering them sequentially through \texttt{ingest}, thereby emulating realistic streaming conditions in which information arrives gradually rather than as a single batch. After ingestion is complete, the agent receives questions that stress four distinct memory operations: \textit{retention} of earlier facts, \textit{integration} of new observations without catastrophic forgetting, \textit{long-range dependency tracking} over distant context segments, and \textit{knowledge revision} through overriding or deleting outdated or conflicting information. Because executing this standardized workflow at scale is computationally expensive, involving large numbers of LLM calls for ingestion, querying, and evaluation, we complement the unified interface with cost-aware execution controls and backend routing mechanisms described above. Together, these efforts ensure controlled, reproducible, and financially tractable evaluation while capturing the full spectrum of memory-centric CRUD-like behaviors, namely the creation of memory traces during ingestion, reading during question answering, updating when new facts appear, and deleting or overriding inconsistent knowledge through selective forgetting.



% Together, these ten subtasks define a unified benchmark suite that captures the multi-faceted nature of memory in LLM-based agents. By spanning recall, incremental learning, long-range reasoning, and controlled memory revision, MemoryAgentBench enables a structured characterization of how agents process, store, and transform information across extended interactions. The benchmark’s competency-driven organization further supports clean comparisons across systems by isolating the specific memory phenomena being stressed, ensuring that variations in behavior can be attributed to memory mechanisms rather than unrelated reasoning differences.




% \begin{table}[t]\centering
% \caption{Statistics of MemoryAgentBench splits with question distributions and truncated context sizes (in thousands of words).}
% \label{tab:memory-bench}
% \footnotesize
% \begin{tabular}{l r r | rrrr | rrrr}
% \toprule
%  & & & \multicolumn{4}{c|}{\textbf{Questions}} & \multicolumn{4}{c}{\textbf{Context (k words)}} \\
% Split & Sessions & Subtasks & Min & Max & Avg & Total & Min & Max & Avg & Total \\
% \midrule
% AR & 22 & 4 & 60 & 100 & 90.9 & 2,000 & 50 & 560 & 206 & 4,551 \\
% TTL & 6 & 2 & 100 & 200 & 116.7 & 700 & 69 & 1,040 & 236 & 1,417 \\
% LRU & 110 & 2 & 1 & 10 & 1.6 & 171 & 48 & 560 & 129 & 14,281 \\
% SF & 8 & 2 & 100 & 100 & 100.0 & 800 & 4 & 183 & 63 & 511 \\
% \bottomrule
% \end{tabular}
% \end{table}


% The statistics in Table \ref{tab:memory-bench} provide a structural summary of MemoryAgentBench and highlight the substantial diversity in session lengths, question density, and context scale across its four competency-aligned splits. Although each split targets a distinct memory phenomenon, the quantitative profile reveals how differently they stress an agent’s memory pipeline during both ingestion and query resolution.

% The \emph{AR} split is relatively compact in the number of sessions (22) but dense in supervision, with roughly 2,000 total questions and context documents ranging from 50k to more than 500k words (median-scale long-context retrieval). Because AR subtasks emphasize precise factual recall and multi-hop retrieval, each question depends on fine-grained memory traces distributed across these high-variance contexts. The \emph{TTL} split is smaller (6 sessions), but its tasks contain longer question sequences—up to 200 questions per session—paired with moderately large contexts (69k–1,040k words). These properties match TTL’s role as an in-session adaptation benchmark, where each additional question provides new concepts, labels, or user preferences that must be integrated without erasing previously learned material. In contrast, the \emph{LRU} split contains by far the largest contexts—110 sessions with an average of 129k words and a total exceeding 14 million words—yet only 1–10 questions per session. This reflects the goal of LRU subtasks: testing global abstraction, long-range temporal and thematic linking, and summarization-like capabilities rather than intensive per-question retrieval. Finally, the \emph{SF} (Selective Forgetting) split uses shorter contexts (4k–183k words) but consistently includes 100 questions per session, ensuring repeated exposure to conflicting or updated facts. This construction directly stresses memory overwriting, contradiction handling, and controlled revision mechanisms.

% Together, these distributions illustrate that MemoryAgentBench does more than aggregate datasets—it operationalizes four qualitatively distinct memory regimes through quantifiable structural differences. AR stresses high-density retrieval over medium to long contexts; TTL couples long sequential supervision with evolving concepts; LRU concentrates its burden in large-document abstraction rather than question volume; and SF forces repeated updates to a manageable but frequently shifting knowledge state. These properties justify evaluating agents uniformly via the \texttt{reset}–\texttt{ingest}–\texttt{query} interface while allowing each split to probe a different axis of memory-centric behavior. The table also highlights why running the full benchmark is computationally demanding: models must repeatedly ingest multi-hundred-kiloword documents and answer thousands of questions whose prompts include large context windows, making MemoryAgentBench a rigorous stress test for both memory design and system efficiency.





% \paragraph{Computational Cost and Configuration Control.}
% \textbf{Benchmark Cost Profile.} MemoryAgentBench is computationally demanding because it is entirely text-based and relies on large language models at multiple stages of the evaluation pipeline. LLMs are invoked not only for agent execution—covering memory initialization, incremental context ingestion, and memory-conditioned question answering—but also as semantic judges during evaluation, where free-form outputs are scored without access to structured labels. Consequently, a single benchmark run involves a large number of LLM calls per session, making runtime and cost highly sensitive to configuration choices. 
% \textbf{Exposed Experimental Controls.} To ensure reproducibility and enable controlled trade-offs between fidelity and efficiency, we expose all key parameters through a unified configuration layer shared across frameworks. The most influential parameter is the maximum number of sessions evaluated per subtask, which directly determines the total number of ingestion and evaluation calls. Equally important is the context window budget, which governs how many tokens are transmitted to the backend model during ingestion and querying, thereby controlling chunk size, overlap, retrieval limits, and the effective memory footprint available to the agent. These parameters allow systematic scaling of experiments—from lightweight diagnostic runs to full benchmark executions—without modifying agent implementations or benchmark logic. 
% \textbf{Backend Routing and Cost-Aware Deployment.} To further support budget- and infrastructure-aware experimentation, we provide two transparent routing backends: a Groq-based router and a local-LLM router. Both forward OpenAI-compatible API calls to alternative backends while preserving the OpenAI interface, allowing any framework that supports OpenAI-style calls to switch backends without internal code changes or framework-specific adaptations. This routing design enables seamless substitution of accelerated cloud inference or fully local models, ensuring that cost constraints or deployment environments do not affect the validity or comparability of experimental results.









% \vspace{-3mm}
% \subsection{Memory Benchmarks}
% \label{subsec:bench_memory}

% % \noindent\textbf{Motivation and Benchmark Selection.}
% Existing memory benchmarks for LLM-based agents evaluate isolated capabilities such as long-context retrieval, dialogue state tracking, preference learning, or counterfactual revision~\cite{hsieh2024ruler,wu2024longmemeval,zhang2024bench,li2018towards,liu2021benchmarking,casanueva2020efficient,larson2019evaluation,he2023large,zhong2023mquake}. While valuable individually, they do not provide a unified view of how agents retain, integrate, generalize, and revise information across extended interactions. We therefore adopt \emph{MemoryAgentBench}~\cite{hu2025memoryagentbench}, which consolidates these dimensions into four core competencies: \textit{Accurate Retrieval (AR)}, \textit{Test-Time Learning (TTL)}, \textit{Long-Range Understanding (LRU)}, and \textit{Selective Forgetting (SF)}.

% \noindent\textbf{Benchmark Structure and Unified Interface.}
% MemoryAgentBench organizes evaluation into four competency-aligned splits composed of adapted subtasks (Table~\ref{tab:memory-bench}). \textit{AR} evaluates factual recall and multi-hop reasoning over long contexts; \textit{TTL} measures in-session learning of concepts and preferences; \textit{LRU} focuses on long-range abstraction and cross-document reasoning; and \textit{SF} probes controlled memory revision via counterfactual updates~\cite{hsieh2024ruler,wu2024longmemeval,zhang2024bench,xu2024detectiveqa,zhong2023mquake}. These splits exhibit substantial variation in session structure, question density, and context length, resulting in distinct memory demands. To enable controlled comparison, the benchmark specifies a standardized agent interface with three methods: \texttt{reset} (memory initialization), \texttt{ingest} (incremental context processing), and \texttt{query} (memory-based question answering), which we implement uniformly across all systems.

% \noindent\textbf{Evaluation Setup.}
% Beyond adopting MemoryAgentBench, MAFBench introduces a unified execution infrastructure that standardizes memory evaluation across agent systems. First, we enforce a common agent interface with \texttt{reset}, \texttt{ingest}, and \texttt{query} operations, enabling consistent session-level execution. Second, we centralize configuration of model parameters, session limits, batching, and scoring to ensure identical experimental conditions and cost control. Third, we replace string-based scoring with LLM-based semantic evaluation to handle diverse answer formats fairly. Fourth, we standardize result logging and aggregation using a shared schema with runtime and token tracking for reproducibility. Finally, we introduce transparent backend routing that redirects compatible API calls to alternative providers, enabling scalable long-context evaluation. In our experiments, memory-intensive requests are routed to Groq-hosted models~\cite{groq} to enable large-scale evaluation at substantially lower cost without modifying framework logic.

\vspace{-1.5mm}
\subsection{Memory Benchmarks}
\label{subsec:bench_memory}

Prior memory benchmarks for LLM-based agents evaluate isolated capabilities, such as long-context retrieval~\cite{hsieh2024ruler,wu2024longmemeval,zhang2024bench}, dialogue and state tracking~\cite{li2018towards,liu2021benchmarking}, incremental preference learning~\cite{casanueva2020efficient,larson2019evaluation,he2023large}, and controlled knowledge revision~\cite{zhong2023mquake}. These benchmarks do not capture how agents retain, integrate, generalize, and revise information across multi-turn interactions. We therefore adopt \emph{MemoryAgentBench}~\cite{hu2025memoryagentbench}, which unifies these dimensions into four core competencies: \textit{Accurate Retrieval (AR)}, \textit{Test-Time Learning (TTL)}, \textit{Long-Range Understanding (LRU)}, and \textit{Selective Forgetting (SF)}.\footnote{\url{https://huggingface.co/datasets/ai-hyz/MemoryAgentBench}}


\noindent\textbf{Benchmark Structure and Unified Interface.}
MemoryAgentBench organizes evaluation into four competency-aligned splits composed of adapted subtasks (Table~\ref{tab:memory-bench}). \textit{AR} measures factual recall and multi-hop reasoning over long contexts; \textit{TTL} evaluates in-session learning of concepts and preferences; \textit{LRU} targets long-range abstraction and cross-document reasoning; and \textit{SF} probes controlled memory revision through counterfactual updates~\cite{hu2025memoryagentbench}.

\noindent\textbf{Evaluation Setup.}
MAFBench integrates MemoryAgentBench into its unified execution pipeline to enable fair and scalable cross-framework comparison. We enforce a standardized agent interface for session-level execution and centralize configuration of model parameters, session limits, batching, and scoring. We also replace string-based metrics with LLM-based semantic evaluation to handle diverse answer formats. Results are logged and aggregated using a shared schema that captures accuracy, runtime, and token usage for reproducibility. To support large-scale long-context evaluation, we introduce transparent backend routing that redirects compatible API calls to alternative providers. In our experiments, memory-intensive workloads are routed to Groq-hosted models~\cite{groq}, enabling lower-cost evaluation without modifying framework implementations.
